\documentclass[10pt,a4paper]{amsart}
\usepackage[margin=19mm]{geometry}
\usepackage{amsmath,amssymb,mathtools}
\usepackage[scr=rsfs,cal=euler]{mathalfa}
\usepackage{xfrac}
\usepackage[unicode,hidelinks,bookmarks=false]{hyperref}
\newcommand{\ie}{i.e.\ }
\newcommand{\cf}{cf.\ }
\newcommand{\resp}{respectively\ }
\newcommand{\Subsection}{\subsection}
\providecommand{\nobreakdash}{\nobreak\hbox{-}}
\setlength{\emergencystretch}{2em}
\multlinegap0pt
\usepackage{amssymb}
\usepackage{mathtools}
\usepackage[shortlabels]{enumitem}
\usepackage{xcolor}
\newtheorem{lemma}{Lemma}
\newtheorem{proposition}{Proposition}
\usepackage{cleveref}
\crefname{lemma}{Lemma}{Lemmas}
\crefname{proposition}{Proposition}{Propositions}
\newcommand{\Fq}{\mathbb F_q}
\newcommand{\Torus}{\mathbb T_q}
\newcommand{\cF}{\mathcal F}
  \def\mathbb#1{#1}%
  \def\mathcal#1{#1}%
\newcommand{\norm}[1]{\lVert#1\rVert}
  \def\operatorname#1{#1}%
\newcommand{\set}[1]{\{#1\}}
\newcommand{\spanq}{\operatorname{span}_{\mathbb F_q}}
\title{Lonely Runners over Function Fields: Quantized Phase--Riesz product}
\author{Xiyu Hu}
\date{Source: arXiv:2608.18818v2 (CC BY 4.0)}
\begin{document}
\maketitle

\noindent\emph{Source and licence.} Lonely Runners over Function Fields: Quantized Phase--Riesz product. Version arXiv:2608.18818v2 (CC BY 4.0), distributed by arXiv under the Creative Commons Attribution 4.0 International licence, http://creativecommons.org/licenses/by/4.0/, as recorded in the arXiv licence metadata for this exact version and shown on the abstract page https://arxiv.org/abs/2608.18818v2. Full licence notice: \texttt{licenses/arxiv-2608.18818-CC-BY-4.0.txt}.\par\smallskip

\noindent\textit{Editorial scope.} Counted unit: the article's proposition prop:compression (source0.tex lines 1121--1134). The article's complete original proof of it is delivered with it (source0.tex lines 1136--1207, one \texttt{proof} environment of 72 lines). No mathematics is written, repaired or re-derived: the statement and the proof are the article's own lines, in the article's own notation, and no retained line is substituted or re-flowed.  Retained uncounted as the source-local context the proof needs: lines 1102--1108.  Nothing was omitted from any retained span, so the package carries no omission record.  No retained line contains a citation, so the wrapper carries no bibliography and no bibliography entry.  No retained line contains a figure environment, an image-inclusion command or drawing code, so no image is delivered and the package declares no asset dependency.  Editorial formatting only, in addition: an AMS \texttt{amsart} preamble that restates the article's own declarations and macros used by the retained lines, namely the environments \texttt{lemma}, \texttt{proposition} on separate counters, together with the standard packages those lines need.  The retained environments print in this document as Lemma 1, Proposition 1 in order of appearance, whereas the article prints the same items with the numbering of its own full document.  The complete native source of the article as distributed in the arXiv e-print for this exact version is bundled unchanged as \texttt{sources/208028/source0.tex}.
\begin{lemma}\label{lem:rational-approx}
Let $P\in\Fq[T]$ be monic and irreducible of degree $M$.  For every
$\alpha\in\Torus$, there is a polynomial $a$ of degree less than $M$ such that
\begin{equation}\label{eq:rational-approx}
 \norm{\alpha-a/P}\le q^{-(M+1)}.
\end{equation}
\end{lemma}

\begin{proposition}\label{prop:compression}
Fix $k\ge2$.  Let $\cF$ be a family of distinct monic polynomials, set
\[
 d:=d(\cF),
 \qquad
 D:=\max_{f\in\cF}\deg f,
\]
and suppose that $\delta(\cF)<q^{-k}$.  Then there is a family $\cF'$ of the same cardinality and the same linear rank such that
\begin{equation}\label{eq:compression-conclusion}
 \delta(\cF')<q^{-k},
 \qquad
 \max_{f'\in\cF'}\deg f'<(k-1)d.
\end{equation}
\end{proposition}

\begin{proof}
We first give one degree-reduction step.  Suppose that
\begin{equation}\label{eq:h-compression}
 h:=\left\lfloor\frac Dd\right\rfloor\ge k-1.
\end{equation}
Put $M=D+1$, choose a monic irreducible polynomial $P$ of degree $M$, and work in the field
\[
 E:=\Fq[T]/(P)\cong\mathbb F_{q^M}.
\]
Let $b_1,\ldots,b_d$ be a basis of
$W:=\spanq\cF$.  Set
\[
 L:=M-h,
 \qquad
 V_L:=\set{g\bmod P:\deg g<L}\le E.
\]
The space $V_L$ has codimension $h$ in $E$.  For each $i$, multiplication by the nonzero residue class $b_i$ is an automorphism of $E$, so
$b_i^{-1}V_L$ also has codimension $h$.  Hence
\begin{equation}\label{eq:compression-intersection}
 \dim_{\Fq}\bigcap_{i=1}^db_i^{-1}V_L
 \ge M-dh\ge1.
\end{equation}
Choose a nonzero residue class $\lambda$ in this intersection.  Since every $f\in W$ is a linear combination of the $b_i$, the residue
$\lambda f\bmod P$ belongs to $V_L$.  Let $g_f$ be its unique representative of degree less than $L$.  Multiplication by $\lambda$ is injective on $W$, so the polynomials $g_f$ are distinct and their span has dimension $d$.  Moreover, two of the $g_f$ cannot be nontrivial scalar multiples: otherwise $g_f=cg_{f'}$ would imply $f=cf'$ because all degrees are less than $M$, and the monicity of $f,f'$ would force $c=1$ and $f=f'$.  Normalising each nonzero $g_f$ by a scalar in $\Fq^\times$ therefore gives a distinct monic family $\cF_1$ of the same size and rank, with
\begin{equation}\label{eq:one-step-degree}
 \max_{g\in\cF_1}\deg g\le D-h.
\end{equation}

It remains to show that the covering property is preserved.  Suppose, to the contrary, that there is an
$\alpha'\in\Torus$ satisfying
\begin{equation}\label{eq:good-time-compressed}
 \norm{\alpha'g}\ge q^{-k}
 \qquad(g\in\cF_1).
\end{equation}
By \cref{lem:rational-approx}, choose $a$ of degree less than $M$ with
\[
 \norm{\alpha'-a/P}\le q^{-(M+1)}.
\]
For $g\in\cF_1$, \eqref{eq:one-step-degree} gives
$\deg g\le M-h-1$, and therefore
\[
 \norm{(\alpha'-a/P)g}
 \le q^{-(h+2)}<q^{-k}.
\]
The ultrametric inequality and \eqref{eq:good-time-compressed} imply
\begin{equation}\label{eq:rational-time-good}
 \norm{(a/P)g}=\norm{\alpha'g}\ge q^{-k}.
\end{equation}

Choose a polynomial representative of $\lambda$.  For each original speed $f$, the corresponding normalised polynomial satisfies
\[
 g_f\equiv c_f\lambda f\pmod P
 \qquad(c_f\in\Fq^\times).
\]
It follows that
\[
 \frac aP g_f
 \equiv
 c_f\frac{a\lambda}{P}f
 \pmod{\Fq[T]}.
\]
Let $\alpha$ be the fractional part of $a\lambda/P$.  Since scalar multiplication by $c_f$ does not change the norm, \eqref{eq:rational-time-good} gives
\[
 \norm{\alpha f}\ge q^{-k}
 \qquad(f\in\cF),
\]
contradicting $\delta(\cF)<q^{-k}$.  Thus
$\delta(\cF_1)<q^{-k}$.

Repeat the reduction while \eqref{eq:h-compression} holds.  At termination,
$\lfloor D/d\rfloor<k-1$, and hence $D<(k-1)d$.  The size and rank are preserved at every step, proving the proposition.
\end{proof}

\end{document}
